%% kapa: exam=1 round=36 year=2025 period=1 form=A subject=economics subject_ko=경제학원론 seq=28 number=68
%% source: Q-Net 감정평가사 기출문제 (cst00309) artlSeq=5231525 file=raw/1008/5231525/2025년 제36회 감정평가사 1차시험 1교시 문제지 원본.pdf
%% conversion: kordoc 4.24.51 -> tools/gen1.py (auto); review=proofread-sonnet; flags=none
\begin{question}{68}
스태그플레이션의 원인과 총수요(AD)-총공급(AS)곡선의 변화를 설명한 것으로 옳지 않은 것을 모두 고른 것은? (단, 폐쇄경제이고, AD곡선은 우하향하고 AS곡선은 우상향한다)
\begin{framedtext}
\begin{tabular}{@{}p{0.34\linewidth}c@{\hspace{1em}}l@{}}
\hspace{2em}\uline{원인} & & \uline{AD-AS곡선변화} \\[0.3em]
ㄱ. 생산비용 상승 & $\Rightarrow$ & AS곡선 왼쪽 이동 \\
ㄴ. 생산비용 하락 & $\Rightarrow$ & AD곡선 오른쪽 이동 \\
ㄷ. 부(-)의 공급충격 & $\Rightarrow$ & AS곡선 왼쪽 이동 \\
ㄹ. 부(-)의 수요충격 & $\Rightarrow$ & AD곡선 오른쪽 이동
\end{tabular}
\end{framedtext}
\begin{choices}
  \item ㄱ, ㄴ
  \item ㄱ, ㄷ
  \item ㄴ, ㄷ
  \item ㄴ, ㄹ
  \item ㄷ, ㄹ
\end{choices}
\end{question}
\answer{4}
